\documentclass[11pt]{article}
\usepackage[margin=1in]{geometry}
\usepackage{amsmath,amssymb}
\usepackage{xcolor}
\usepackage[colorlinks=true,linkcolor=blue!50!black,citecolor=blue!50!black,urlcolor=blue!50!black]{hyperref}

\newcommand{\F}{\mathbb{F}}
\newcommand{\eq}{\mathrm{eq}}
\newcommand{\pkd}{\mathrm{packed}}

\title{\textbf{Ring switching, generalized}}

\author{}
\date{}

\begin{document}
\maketitle
\vspace{-3em}
\begin{center}
How to transform a PCS for $\F_{2^f}$ into a PCS for $\F_{2^e}$, for arbitrary $e$ and $f$.
\end{center}

\paragraph{Goal.} We want $f$ \emph{bit}-valued multilinears $P_1,\dots,P_f:\{0,1\}^m\to\{0,1\}$ over the same $m$ variables, to be committed through a dense PCS over the \emph{packing} field $\F_{2^f}$ (if we can commit to bits, we can trivially commit in any extension). The evaluation claims, however, arrive over a \emph{different} field: at an opening point $r\in\F_{2^e}^m$ we are given $P_i(r)=\alpha_i\in\F_{2^e}$. Ring switching~\cite{DP26} packs the $f$ polynomials into one $\F_{2^f}$-dense polynomial and reduces the $f$ claims to a single claim on it. The two extension degrees, $f$ (packing) and $e$ (opening), are unrelated (not necessarily a sub-field / extension-field).

\paragraph{Packing.} Fix an $\F_2$-basis $\{b^{F}_i\}_{i=1}^f$ of $\F_{2^f}$. Consider the multilinear polynomial:
\[
P_\pkd(y)=\sum_{i=1}^f P_i(y)\,b^{F}_i,\qquad y\in\{0,1\}^m,
\]
dense over $\F_{2^f}$, which we commit to.

\paragraph{Reduction.} The claims share a \emph{common} point $r\in\F_{2^e}^m$; each is
\[
\alpha_i=P_i(r)=\sum_{y\in\{0,1\}^m}\eq(r,y)\,P_i(y),
\]
with bits $P_i(y)$ and weights $\eq(r,y)\in\F_{2^e}$. Fix an $\F_2$-basis $\{b^{E}_w\}_{w=1}^e$ of the \emph{opening} field $\F_{2^e}$ and expand the weights, $\eq(r,y)=\sum_{w=1}^e A(y,w)\,b^{E}_w$, with $A:\{0,1\}^m\times[e]\to\F_2$ public, multilinear in $y$ for each fixed $w$, with efficiently computable MLE (see \nameref{section:eval_B} below). Substituting and pulling out the basis,
\[
\alpha_i=\sum_{w=1}^eb^{E}_w\underbrace{\sum_{y\in\{0,1\}^m}A(y,w)\,P_i(y)}_{\in\,\F_2},
\]
so the $w$-th $\F_2$-coordinate of $\alpha_i$ is $\alpha_{w,i}=\sum_{y\in\{0,1\}^m}A(y,w)\,P_i(y)$. Now recombine \emph{across $i$} with the \emph{packing} basis $\{b^{F}_i\}$: define, for each $w \in \{1,\dots,e\}$,
\[
t_w:=\sum_{i=1}^f\alpha_{w,i}\,b^{F}_i=\sum_{y\in\{0,1\}^m}A(y,w)\underbrace{\sum_{i=1}^fP_i(y)\,b^{F}_i}_{=\,P_\pkd(y)}=\sum_{y\in\{0,1\}^m}A(y,w)\,P_\pkd(y)\ \in\ \F_{2^f}
\]
one claim per $w$, now purely about the packed polynomial $P_\pkd$ over $\F_{2^f}$ and public $A$. Batch these $e$ claims with powers of a random $\gamma\leftarrow\F_{2^f}$ into a single sumcheck on $\sum_{y\in\{0,1\}^m}B(y)\,P_\pkd(y)=\sum_{w=1}^e\gamma^w t_w$, with $B(y)=\sum_{w=1}^e\gamma^w A(y,w)$. Its output is the packed claim $P_\pkd(r')=\alpha'$ at a point $r'\in\F_{2^f}^m$ (plus a $B(r')$ the verifier computes itself). Done.

\paragraph{Soundness.} Suppose the input is wrong: $\alpha_i\neq P_i(r)$ for some $i$, against the committed $P_\pkd$. Then the reduction outputs a correct packed claim $P_\pkd(r')=\alpha'$ with probability at most
\[
\underbrace{\frac{e}{2^f}}_{\text{batching}}\;+\;\underbrace{\frac{2m}{2^f}}_{\text{sumcheck}}\;=\;\frac{2m+e}{2^f},
\]

Proof:

\begin{enumerate}
\item The map $(\alpha_i)_{i=1}^f\mapsto(t_w)_{w=1}^e$ is a \emph{bijection} $\F_{2^e}^f\to\F_{2^f}^e$. So a wrong input forces, with certainty, $t_w\neq\sum_{y\in\{0,1\}^m}A(y,w)\,P_\pkd(y)$ for some $w$.
\item By Schwartz--Zippel, $\sum_{w=1}^e\gamma^w t_w$ then differs from $\sum_{y\in\{0,1\}^m}B(y)\,P_\pkd(y)$ except with probability $\le e/2^f$ over $\gamma\leftarrow\F_{2^f}$.
\item Sumcheck started from this false value (summand of degree $2$, over $m$ variables) ends at a correct claim $P_\pkd(r')=\alpha'$ with probability $\le 2m/2^f$.
\end{enumerate}

\phantomsection
\section*{Evaluating $B$}
\label{section:eval_B}

It remains to compute $B(r')$, where $B(y)=\sum_{w=1}^e\gamma^w A(y,w)$ and $A(y,w)\in\F_2$ is the $w$-th coordinate of $\eq(r,y)$, i.e.\ $\eq(r,y)=\sum_{w=1}^e A(y,w)\,b^E_w$. Following~\cite{Flock}, this gives $B$ a closed form on the hypercube. Let $\Phi:\F_{2^e}\to\F_{2^f}$ be the $\F_2$-linear map with $\Phi(b^E_w)=\gamma^w$. For every $y\in\{0,1\}^m$:
\[
\Phi\big(\eq(r,y)\big)=\Phi\Big(\sum_{w=1}^e A(y,w)\,b^E_w\Big)=\sum_{w=1}^e A(y,w)\,\Phi(b^E_w)=\sum_{w=1}^e A(y,w)\,\gamma^w=B(y).
\]
So on the hypercube $B$ is ``compute $\eq(r,y)\in\F_{2^e}$, then apply $\Phi$''.

\paragraph{Closed form.} Two elementary facts do all the work.

\emph{Fact 1 ($\eq$ in characteristic $2$).} In each factor of $\eq$ the two $x_jy_j$ terms cancel:
\[
\eq(x,y)=\prod_{j=1}^m\bigl(x_jy_j+(1+x_j)(1+y_j)\bigr)=\prod_{j=1}^m(1+x_j+y_j).
\]

\emph{Fact 2 (linear maps are sums of squarings).} View $\Phi$ as mapping into $L=\F_{2^{\mathrm{lcm}(e,f)}}$, the smallest field containing both $\F_{2^e}$ and $\F_{2^f}$ (for $e=f$: $L=\F_{2^e}$, nothing to embed). Squaring is additive in characteristic $2$, so any \emph{linearized polynomial}~\cite{WikiLinPoly}
\[
a\;\mapsto\;c_0\,a+c_1\,a^2+c_2\,a^4+\dots+c_{e-1}\,a^{2^{e-1}},\qquad c_k\in L,
\]
is an $\F_2$-linear map $\F_{2^e}\to L$. We claim $\Phi$ is of this form, and find the $c_k$ by writing out what $\Phi$ must do on the basis, $\Phi(b^E_w)=\gamma^w$:
\[
\sum_{k=0}^{e-1}\bigl(b^E_w\bigr)^{2^k}c_k=\gamma^w,\qquad w=1,\dots,e:
\]
$e$ linear equations in the $e$ unknowns $c_k$. The coefficient matrix $\bigl((b^E_w)^{2^k}\bigr)_{w,k}$ is a \emph{Moore matrix}, invertible exactly because the $b^E_w$ are $\F_2$-linearly independent~\cite{WikiMoore}, so the $c_k$ exist and are unique; the resulting linearized polynomial agrees with $\Phi$ on an $\F_2$-basis, hence (both being $\F_2$-linear) everywhere. The matrix depends only on the basis, not on $\gamma$: invert it once and for all; each reduction then reads off its $(c_k)_{k<e}$ by one matrix-vector product, $O(e^2)$ multiplications.

Now compute. Squaring respects sums and products and fixes bits ($y_j^2=y_j$), so on the hypercube, by Fact~1,
\[
\eq(r,y)^{2^k}=\prod_{j=1}^m(1+r_j+y_j)^{2^k}=\prod_{j=1}^m\bigl(1+r_j^{2^k}+y_j\bigr)=\eq\bigl(r^{2^k},y\bigr),
\qquad r^{2^k}:=\bigl(r_1^{2^k},\dots,r_m^{2^k}\bigr).
\]
Since $B$ is multilinear, evaluating it at $r'$ and expanding with Fact~2:
\[
B(r')=\sum_{y\in\{0,1\}^m}B(y)\,\eq(r',y)
=\sum_{y\in\{0,1\}^m}\Phi\bigl(\eq(r,y)\bigr)\,\eq(r',y)
=\sum_{k<e}c_k\sum_{y\in\{0,1\}^m}\eq\bigl(r^{2^k},y\bigr)\,\eq(r',y).
\]
The inner sum evaluates, at $r'$, the multilinear extension of $y\mapsto\eq(r^{2^k},y)$, which is $\eq(r^{2^k},\cdot\,)$ itself. So it collapses to $\eq(r^{2^k},r')$, and by Fact~1 again:
\[
\boxed{\ \
B(r') \;=\; \sum_{k<e}\, c_k \prod_{j=1}^m
\bigl( 1 + r_j^{2^k} + r'_j \bigr).
\ \ }
\]

\paragraph{Cost.} Level $k$ needs one squaring per coordinate (to get $r_j^{2^k}$ from $r_j^{2^{k-1}}$) and one $m$-factor product: $O(em)$ multiplications in $L$, on top of the $O(e^2)$ for the $c_k$: plain field arithmetic throughout, no bit decompositions.

\section*{Assisted recursive verification}

The linearized-polynomial computation above can be expensive inside a recursive
verifier.  Let $(B_j)_{j<f}$ be the packing basis and
$(\delta_i)_{i<f}$ its trace-dual basis:
\[
  \operatorname{Tr}(\delta_i B_j)=[i=j].
\]
For a ring-switch weight vector $w_t=(w_t(i))_{i<f}$, define
\[
  c_t(k)=\sum_{i<f}w_t(i)\,\delta_i^{2^k},
  \qquad
  L_t(y)=\sum_{k<f}c_t(k)y^{2^k}.
\]
Constructing the whole vector $c_t$ directly costs $\Theta(f^2)$ field
operations.

\paragraph{Certifying assisted coefficients.}
The coefficients may instead be supplied as advice and certified through the
fixed table
\[
  D(k,i)=\delta_i^{2^k}.
\]
Let $d=\lceil\log_2 f\rceil$ and pad $D$ and every length-$f$ vector with zeros
to a $2^d$-by-$2^d$ table.  After observing all advised vectors, sample
batching scalars $\gamma_t$ and a random row point $\rho\in\F_{2^f}^d$.  When
$w_t(i)=\eq(r''_t,i)$, the identity
\[
  \sum_t\gamma_t\sum_k\eq(\rho,k)c_t(k)
  =
  \sum_i \widetilde D(\rho,i)
       \sum_t\gamma_t\eq(r''_t,i)
\]
is checked by a $d$-round product sumcheck over $i$.  It leaves one evaluation
of the fixed $2d$-variable polynomial $\widetilde D$, which may be evaluated
outside the recursive circuit.  Thus arbitrarily many coefficient vectors
reduce to one constant-size fixed-table claim.

\paragraph{Proof-dependent relations.}
The same coefficients often appear in two further verifier computations.  If
$s_{t,a}=(s_{t,a,j})_{j<f}$ is a proof message, its transpose contribution is
\[
  T_{t,a}
  =
  \sum_{j<f}B_j\,L_t(s_{t,a,j})
  =
  \sum_{j,k<f}B_j\,c_t(k)\,s_{t,a,j}^{2^k}.
\]
If $q_t=(q_{t,j})_{j<\ell_t}$ is the final PCS query, the corresponding
terminal weight has the form
\[
  E_{t,a}
  =
  \sum_{k<f}c_t(k)
  \prod_{j<\ell_t}\bigl(z_{t,a,j}^{2^k}+1+q_{t,j}\bigr).
\]
These identities cost $\Theta(f^2)$ and $\Theta(f\ell_t)$ respectively, but
unlike the coefficient transform their right-hand sides depend on proof
messages.  An assisted verifier may use advised values $T_{t,a}$ and
$E_{t,a}$, bind all their inputs into the recursive statement, and let an
outer verifier recompute them.

\paragraph{Composability tradeoff.}
The fixed-table reduction is naturally batchable and remains constant-size
across recursion layers.  Forwarding the proof-dependent inputs above instead
adds one record per immediate proof.  Those records cannot be compressed merely
by random-linear combination, because their right-hand sides are nonlinear and
different for every proof.  Constant-size deep recursion therefore requires
either checking these identities inside the recursive circuit or giving them a
separate proof-carrying accumulator.  In practice, assistance is worthwhile
only when the saved recursive work justifies this extra protocol and statement
complexity.

\section*{Remarks}

\begin{enumerate}
\item \emph{Small $\F_{2^f}$.} If $2^f$ is too small for the error $(2m+e)/2^f$, sample $\gamma$ (and run the sumcheck) in an extension $L\supseteq\F_{2^f}$: the error becomes $(2m+e)/|L|$, and $P_\pkd$ is just opened at $r'\in L^m$.
\item \emph{Any field.} Nothing uses characteristic $2$: replace $\F_2$ by any field $K$ and $\F_{2^f},\F_{2^e}$ by degree-$f$ and degree-$e$ extensions of $K$; the reduction still works, and efficient algorithms for evaluating $B$ exist in general too.
\end{enumerate}

\section*{Credits}

The idea of generalizing Ring-Switching, originally described in \cite{DP26}, to arbitrary extension degrees (non necessarily powers of two), comes from Lev Soukhanov, \textit{[[alloc]init]}.

\begin{thebibliography}{9}
\bibitem{DP26} Benjamin E. Diamond and Jim Posen. Polylogarithmic proofs for multilinears over binary towers. In \emph{Advances in Cryptology -- EUROCRYPT 2026}, pages 3--32. Springer, 2026.
\bibitem{Flock} Benedikt B\"unz, Ron Rothblum, and William Wang. Flock: Fast proving for batch Boolean computations. Cryptology ePrint Archive, Paper 2026/1329, 2026. \url{https://eprint.iacr.org/2026/1329}.
\bibitem{WikiLinPoly} Wikipedia. Linearised polynomial. \url{https://en.wikipedia.org/wiki/Linearised_polynomial}.
\bibitem{WikiMoore} Wikipedia. Moore matrix. \url{https://en.wikipedia.org/wiki/Moore_matrix}.
\end{thebibliography}

\end{document}
